\section{Compact plateaux and strict increase}\label{sec:trapping}

We now apply Corollary~\ref{cor:activity} to the reaction network.  For each
sufficiently small scale we will construct a compact set in the positive
orthant that is forward invariant.  For a fixed positive stoichiometric
class, we will also prove that a trajectory in that class remaining in
the scale's box reaches this set in finite time.  The affine
minimum and sign argument adapt the trapping construction of
\cite[Section~3]{OpenAI2026}; the strict increase away from a plateau is the
additional property needed for finite entry within a fixed class.

Use $X_h$ from Section~\ref{sec:affine} with $m=d$, and for $0<h<1$ put
\[
 B_h=[h,h^{-1}]^d.
\]
Every $x\in B_h$ is uniquely $X_h(p)$ for $p\in[-1,1]^d$, with
$p=\log x/\log h$, where the logarithm of a vector is taken coordinatewise.
We first choose the tolerance in Corollary~\ref{cor:activity} so that active
slopes determine both coordinate signs and the ordering of monomials.
Let
\begin{equation}\label{eq:comparison-directions}
 D=\{e_i:1\leq i\leq d\}
   \cup\{z-y:y,z\in C,\ z\ne y\},
\end{equation}
where $e_i$ is the $i$th coordinate vector.  This is a finite set of nonzero
vectors.  We include all pairs of complexes because the flux estimate
will compare sources of different reactions, possibly in different
components of the reaction graph.  Define, for $r\in[-1,1]^d$,
\begin{equation}\label{eq:network-tolerance}
 E(r)=\min\left(
 \left\{\frac12\right\}\cup
 \left\{\frac{\abs{r\cdot v}}{2\norm{v}_1}:
                 v\in D,\ r\cdot v\ne0\right\}\right).
\end{equation}
Zero scalar products impose no sign restriction and are omitted.
Thus $E$ is positive at every point, including on the comparison
hyperplanes; no continuity is needed.
Fix the finite nonempty family supplied by
Corollary~\ref{cor:activity} for this $E$, and write its affine functions as
\[
 L_{\lambda,h}(x)=r_\lambda\cdot x+b_\lambda(h),
 \qquad
 F_h(x)=\min_{\lambda\in\Lambda}L_{\lambda,h}(x).
\]
The slopes $r_\lambda\in[-1,1]^d$ and the index set $\Lambda$ are fixed
independently of $h$.  A label $\lambda$ is active at $x$ if
$L_{\lambda,h}(x)=F_h(x)$; denote the set of these labels by $A_h(x)$.
For every sufficiently small $h$, the activity estimate is
\begin{equation}\label{eq:network-activity}
 x=X_h(p)\in B_h,\quad \lambda\in A_h(x)
 \quad\Longrightarrow\quad
 \norm{p-r_\lambda}_\infty<E(r_\lambda).
\end{equation}

Because the family of slopes is fixed and finite, the number
\begin{equation}\label{eq:uniform-sign-gap}
 H=\min\left(
 \{1\}\cup
 \left\{\frac{\abs{r_\lambda\cdot v}}2:
       \lambda\in\Lambda,\ v\in D,\ r_\lambda\cdot v\ne0\right\}\right)
\end{equation}
is positive.  If $x=X_h(p)\in B_h$, $\lambda\in A_h(x)$, and
$r=r_\lambda$, then
\[
 \abs{(p-r)\cdot v}
 \leq\norm{p-r}_\infty\norm{v}_1
 <E(r)\norm{v}_1
 \leq\frac{\abs{r\cdot v}}2
 \qquad (v\in D,\ r\cdot v\ne0).
\]
Consequently,
\begin{equation}\label{eq:sign-and-size}
 r\cdot v\ne0
 \quad\Longrightarrow\quad
 \operatorname{sgn}(p\cdot v)=\operatorname{sgn}(r\cdot v),
 \qquad \abs{p\cdot v}\geq H
 \qquad (v\in D).
\end{equation}

Fix a positive stoichiometric class $P$, and choose one representative
$c\in P$ once for this class.  Choose $h_P\in(0,1)$ so small that the
affine functions are defined and \eqref{eq:network-activity} holds for
every $0<h\leq h_P$, and also
\begin{equation}\label{eq:class-scale}
 \norm{\frac{\log c}{\log h}}_\infty<H,
 \qquad
 h^H\sum_{e\in R}k_e<\min_{e\in R}k_e
 \quad (0<h\leq h_P).
\end{equation}
The second inequality is imposed only when $R\ne\varnothing$.
Such a choice exists because $H>0$, all the rates are positive, and
both $1/\abs{\log h}$ and $h^H$ tend to zero as $h$ tends to zero.
In particular, $c\in\operatorname{int}B_h$ for these scales.
All the remaining statements of this section concern $0<h\leq h_P$.

\begin{proposition}[The compact plateau]\label{prop:plateau}
For $0<h\leq h_P$, let
\[
 M_h=F_h(c),
 \qquad
 K_h=\{x\in\R^d:F_h(x)=M_h\}.
\]
There is a constant label with value $M_h$, and $F_h\leq M_h$ on
$\R^d$.  The set $K_h$ is nonempty, compact, and convex, and
\begin{equation}\label{eq:plateau-containment}
 K_h\subset\operatorname{int}B_{\sqrt h}.
\end{equation}
\end{proposition}

\begin{proof}
Write $p_c=\log c/\log h$.  If a label active at $c$ had a nonzero
coordinate $r_{\lambda,i}$, \eqref{eq:sign-and-size} with $v=e_i$ would
give $\abs{p_{c,i}}\geq H$, contrary to \eqref{eq:class-scale}.
Thus every label active at $c$ has zero slope.  At least one label is
active there, so one of the affine functions is the constant $M_h$.
It follows that $F_h\leq M_h$ everywhere, and that $M_h$ is the least
value of any constant label.  Hence
\[
 K_h=\bigcap_{\lambda\in\Lambda}
       \{x\in\R^d:L_{\lambda,h}(x)\geq M_h\}
\]
is closed and convex, and contains $c$.

If $x=X_h(p)\in K_h\cap B_h$, the constant label with value $M_h$ is
active at $x$.  Applying \eqref{eq:network-activity} to its zero slope
gives
\[
 \norm{p}_\infty<E(0)=\frac12,
 \qquad\text{and thus}\qquad
 x\in\operatorname{int}B_{\sqrt h}.
\]
In particular, $K_h$ misses the boundary of $B_h$.  If some point of
$K_h$ lay outside $B_h$, the segment joining it to
$c\in\operatorname{int}B_h$ would meet that boundary at a point of
$K_h$, by convexity.  This is impossible.  Thus $K_h\subset B_h$, and
the preceding argument proves \eqref{eq:plateau-containment} for every
point of $K_h$.  Closedness and boundedness give compactness.
\end{proof}

The plateau is now a compact candidate for a trapping set. A slope in
$S^\perp$ gives an affine function that is constant on $P$. We first
exclude such slopes below the plateau, then compare reaction fluxes to
prove strict increase for every remaining active label.

\begin{proposition}[Increase of active affine functions]\label{prop:increase}
Fix $0<h\leq h_P$.  For every $x\in B_h$ and every
$\lambda\in A_h(x)$,
\begin{equation}\label{eq:active-flux-sign}
 r_\lambda\cdot f(x)\geq0,
 \qquad
 r_\lambda\cdot f(x)>0\quad\text{if }r_\lambda\notin S^\perp.
\end{equation}
Moreover, there is a number $\delta_h>0$ such that
\begin{equation}\label{eq:uniform-active-increase}
 r_\lambda\cdot f(x)\geq\delta_h
 \quad\text{whenever }x\in B_h,\ \lambda\in A_h(x),
                       \ r_\lambda\notin S^\perp.
\end{equation}
At every point of $P\setminus K_h$, all active slopes lie outside
$S^\perp$.  In particular, \eqref{eq:uniform-active-increase} applies
to every active label at each $x\in B_h\cap(P\setminus K_h)$.
\end{proposition}

\begin{proof}
First, let $x\in P\setminus K_h$.  If $r_\lambda\in S^\perp$,
then $x-c\in S$ implies
\[
 L_{\lambda,h}(x)=L_{\lambda,h}(c)\geq F_h(c)=M_h.
\]
But $F_h(x)<M_h$ by Proposition~\ref{prop:plateau}.  Such a label
cannot be active at $x$.

If $R$ is empty, then $f=0$ and $S=\{0\}$, so
\eqref{eq:active-flux-sign} holds and
\eqref{eq:uniform-active-increase} is vacuous.  For the flux comparison,
assume $R\ne\varnothing$ and fix $x=X_h(p)\in B_h$ and an active slope $r$.
Whenever $y,z\in C$ satisfy $r\cdot y<r\cdot z$,
\eqref{eq:sign-and-size} applied to $z-y$ gives
\begin{equation}\label{eq:monomial-order}
 \frac{x^z}{x^y}=h^{p\cdot(z-y)}\leq h^H.
\end{equation}
List the distinct values of $r\cdot y$, for $y\in C$, as
$a_1<\cdots<a_m$.  For $1\leq j<m$, define the net upward flux across
the cut at $a_j$ by
\begin{align*}
 Q_j(x)={}&
 \sum_{\substack{u\to u'\in R\,:\\r\cdot u\leq a_j<r\cdot u'}}
       k_{u\to u'}x^u
 -\sum_{\substack{u\to u'\in R\,:\\r\cdot u'\leq a_j<r\cdot u}}
       k_{u\to u'}x^u.
\end{align*}
An edge crossing the cut in the first sum is called upward, and an
edge in the second sum is called downward.

If there is a downward edge $u\to u'$, weak reversibility supplies a
directed path from $u'$ back to $u$.  This path starts below the cut
and ends above it, so it contains an upward edge, say $w\to w'$.
Every source $u$ of a downward edge has $r\cdot u>r\cdot w$.
By \eqref{eq:monomial-order} and \eqref{eq:class-scale},
\[
 \sum_{\substack{u\to u'\in R\,:\\r\cdot u'\leq a_j<r\cdot u}}
       k_{u\to u'}x^u
 \leq h^H x^w\sum_{e\in R}k_e
 < k_{w\to w'}x^w.
\]
Thus one upward edge already contributes more than the entire
downward sum, and $Q_j(x)>0$.  If there is no downward edge, then
$Q_j(x)\geq0$, with strict inequality whenever there is an upward
edge.  In all cases $Q_j(x)$ is nonnegative, and it is positive if
any reaction crosses the cut.

The contribution of an edge to $r\cdot f(x)$ is its flux multiplied
by the difference between the levels of its target and source.
Writing this difference as the signed sum of the intervening gaps gives
\begin{equation}\label{eq:cut-flux-decomposition}
 r\cdot f(x)=\sum_{j=1}^{m-1}(a_{j+1}-a_j)Q_j(x).
\end{equation}
Every term is nonnegative.  If $r\notin S^\perp$, some reaction
$y\to y'$ satisfies $r\cdot(y'-y)\ne0$ and therefore crosses at least
one cut; the corresponding term is positive.  This proves
\eqref{eq:active-flux-sign}.  When $m=1$, the sum is empty and every
reaction has zero scalar product with $r$, as required.  Self-loops
never cross a cut and contribute zero to the vector field.

For the uniform bound, the active set of a fixed label,
\[
 T_{\lambda,h}=
 \{x\in B_h:L_{\lambda,h}(x)=F_h(x)\},
\]
is compact.  If $r_\lambda\notin S^\perp$ and $T_{\lambda,h}$ is
nonempty, the continuous function $r_\lambda\cdot f$ is strictly
positive on this set and has a positive minimum there.  Take the
minimum of these positive numbers over the finitely many eligible
labels.  If there are none, choose any $\delta_h>0$.  This proves
\eqref{eq:uniform-active-increase}, including the case $R=\varnothing$.
\end{proof}

The strict inequality in Proposition~\ref{prop:increase} is uniform at
a fixed scale even near the plateau.  Together with the finite nature
of the affine minimum, this gives a direct finite-time trapping
statement.

\begin{lemma}[Trapping and finite entry]\label{lem:trapping}
Fix $0<h\leq h_P$.  Along any solution segment contained in $B_h$,
the function $t\mapsto F_h(x(t))$ is nondecreasing.  If the segment
$x([s,t])$ is contained in $B_h\cap(P\setminus K_h)$, then
\begin{equation}\label{eq:trajectory-increase}
 F_h(x(t))-F_h(x(s))\geq\delta_h(t-s).
\end{equation}
Every solution starting in $K_h$ exists for all nonnegative time and
remains in $K_h$.  In particular, solutions starting in $K_h\cap P$
remain in $K_h\cap P$.  Finally, any global solution in $P$ that stays
in $B_h$ from some time onward enters $K_h$ in finite time and then
remains there.
\end{lemma}

\begin{proof}
For a solution $x(t)$ and each fixed time $t$ before its maximal
existence time, the right derivative of its affine minimum is
\begin{equation}\label{eq:right-derivative}
 \frac{d^+}{dt}F_h(x(t))
   =\min_{\lambda\in A_h(x(t))}r_\lambda\cdot f(x(t)).
\end{equation}
Indeed, all inactive labels have a positive gap above the minimum at
time $t$.  There are finitely many of them, so continuity excludes
them from the minimum for a sufficiently short time to the right.
For each active label, differentiability of the solution gives
\[
 L_{\lambda,h}(x(t+s))
 =F_h(x(t))+s\,r_\lambda\cdot f(x(t))+o(s)
 \qquad(s\downarrow0).
\]
Taking the minimum of these finitely many expansions proves
\eqref{eq:right-derivative}.

We use the elementary fact that a continuous function $g$ on
$[a,b]$ with right derivative at least $q$ at every point of $[a,b)$
satisfies $g(b)-g(a)\geq q(b-a)$.  For completeness, for any
$\eta>0$ the function $g(u)-qu+\eta u$ has positive right derivative
before $b$.  Its maximum on the compact interval therefore cannot
occur before $b$.  Comparing its values at $a$ and $b$ and then
letting $\eta\downarrow0$ proves the assertion.  Formula
\eqref{eq:right-derivative} and Proposition~\ref{prop:increase} now
give monotonicity in $B_h$, and give \eqref{eq:trajectory-increase}
when the segment lies in $P\setminus K_h$.

Suppose next that $x(0)\in K_h$, and consider the maximal solution of
the polynomial vector field $f$.  Since
$K_h\subset\operatorname{int}B_h$, the solution initially stays in
$B_h$.  On any interval on which it does so, monotonicity and the
global upper bound $F_h\leq M_h$ force
$F_h(x(t))=M_h$.  Hence it remains in $K_h$.  It cannot have a first
exit from $\operatorname{int}B_h$ within its interval of existence:
continuity would place the exit point in both $K_h$ and the boundary
of $B_h$, contradicting Proposition~\ref{prop:plateau}.
Thus the solution stays in the compact set $K_h$ throughout its
maximal interval.  A solution of a smooth vector field that remains
in a compact subset of its domain extends past any finite endpoint
of its interval of existence.  To see the relevant point here, $f$
is bounded on a neighborhood of $K_h$, so a finite endpoint would
give a limit of $x(t)$ in $K_h$, at which local existence continues
the solution.  The solution is therefore global.  It remains
positive by \eqref{eq:plateau-containment}; moreover, $f(x)\in S$
implies $x(t)-x(0)\in S$, so a solution starting in $P$ stays in $P$.

Finally, let a global solution in $P$ remain in $B_h$ for all
$t\geq t_0$.  If it never enters $K_h$ after $t_0$, then
\eqref{eq:trajectory-increase} gives
\[
 F_h(x(t))\geq F_h(x(t_0))+\delta_h(t-t_0)
 \qquad(t\geq t_0),
\]
contradicting $F_h\leq M_h$ for sufficiently large $t$.
It must enter $K_h$ at a finite time, and the forward invariance just
proved keeps it there.
\end{proof}
